% ---------- ¶6  data and vantage point  [PRESERVED from current ¶6] ----------
Testing a theory of strategic concealment requires observing both sides of the information
asymmetry at once, the seller's true information, the seller's disclosure choice and the
market's response. Most empirical work observes only one or two. Studies of information
asymmetry in housing infer agent advantages from equilibrium price deviations
\citep{levitt_syverson_2008}; analyses of certification intermediaries characterise incentives
without observing individual information endowments \citep{lizzeri_1999}; and work on the
energy-efficiency gap proxies for seller awareness through indirect variation in energy prices
\citep{myers_asymmetric_2020}. We observe all three. Listing records from SeLoger, a major
French residential platform, give the listing price, property characteristics, the seller's
disclosure decision, and, when a score is displayed, the letter itself. The registry of the
\emph{Agence de la Transition \'Ecologique} (ADEME) is the official national database of
certified DPE audits, reported by accredited third-party assessors and compiled independently
of the listing market. We match listings to the registry on address, property type, area, and
floor, which identifies whether a certified audit existed at the time of listing, the one
thing buyers can never observe. Because the match requires a geocoded address, the matched
sample is the subset of listings that can be located at the address level rather than the full
listing universe. Within that matched sample, every listing falls into one of four groups: truthful
disclosers, strategic hiders who hold a certificate and omit it, genuinely uninformed
sellers for whom no certificate exists, and a small residual displaying a score without a
verified audit. The information structure of the disclosure game is therefore observed rather
than inferred, a vantage point that, to our knowledge, is unique in this literature. Transaction
records enter separately and are not matched to individual listings. We use them to estimate
the mapping from listing to realised sale price, which supports treating the listing price as
the relevant outcome, since the negotiation margin in French residential sales is small and
most properties transact at or near list price
\citep{merlo_bargaining_2004, genesove_loss_2001, friggit_marche_2013}.
